% Part: computability
% Chapter: computability-theory
% Section: introduction

\documentclass[../../../include/open-logic-section]{subfiles}

\begin{document}

\olfileid[fr]{cmp}{thy}{int}
\olsection{Introduction}

La branche de la logique appelée \emph{théorie de la calculabilité}
étudie la calculabilité des fonctions et des ensembles, ainsi que
leur calculabilité relative. Témoignage de l'influence de Kleene,
elle portait autrefois le nom de \emph{théorie de la récursion};
les deux appellations restent aujourd'hui en usage.

Appelons une fonction $f\colon\Nat\pto\Nat$
\emph{partielle calculable} lorsqu'elle peut être calculée
dans un modèle usuel de calcul, par exemple celui des machines de Turing. Si $f$ est totale, nous dirons
simplement qu'elle est \emph{calculable}. Une relation~$R$
dont la fonction caractéristique $\Char{R}$ est calculable
est également dite calculable. Pour des fonctions partielles
$f$ et~$g$, nous écrirons $f(x) \fdefined$ si $f$ est
définie en~$x$, autrement dit si $x$ appartient à son
domaine, et $f(x) \fundefined$ dans le cas contraire.
La notation $f(x) \simeq g(x)$ signifie que les deux
valeurs sont indéfinies ou qu'elles sont toutes deux
définies et égales.

On peut développer la théorie de la calculabilité sans
se référer constamment à un modèle de calcul particulier.
À cette fin, on montre qu'il existe une fonction partielle
calculable universelle~$\fn{Un}(k,x)$. Elle permet
d'énumérer les fonctions partielles calculables. Nous
noterons $\cfind{k}$ la $k$-ième fonction partielle
calculable unaire, définie par
$\cfind{k}(x) \simeq \fn{Un}(k,x)$. (Kleene employait
pour cela la notation $\{ k \}$, moins usitée de nos
jours.) Plus généralement, on peut énumérer uniformément
les fonctions partielles calculables de toutes les arités;
$\cfind{k}[n]$ désignera la $k$-ième fonction partielle
récursive à $n$~arguments.

Si $f(\vec x,y)$ est une fonction totale ou partielle,
$\umin{y}{f (\vec x,y)}$ est la fonction de~$\vec x$
qui renvoie le plus petit~$y$ tel que $f(\vec x,y)=0$,
à condition que toutes les valeurs $f(\vec x,0),\dots,
f(\vec x,y-1)$ soient définies. S'il n'existe pas de
tel~$y$, $\umin{y}{f (\vec x,y)}$ est indéfini.
Pour une relation $R(\vec x,y)$, la notation
$\umin{y}{R(\vec x,y)}$ désigne le plus petit~$y$
tel que $R(\vec x,y)$ soit vraie, c'est-à-dire le
plus petit~$y$ tel que
$1 \tsub \Char{R}(\vec x,y)=0$.

Pour montrer qu'une fonction est calculable, on peut
procéder de deux manières :
\begin{enumerate}
\item De façon formelle : décrire explicitement une
  machine de Turing ou une fonction récursive partielle,
  puis montrer qu'elle calcule la fonction voulue.
\item De façon informelle : décrire un algorithme qui
  la calcule et invoquer la thèse de Church.
\end{enumerate}
La frontière entre ces deux démarches n'est pas nette.
Une description assez détaillée de l'algorithme doit
permettre de voir comment construire, en principe, la
machine de Turing ou la suite de définitions récursives
partielles appropriée. Des définitions entièrement formelles
seraient rarement éclairantes; nous chercherons donc un
juste milieu, en donnant des preuves intuitives mais
assez précises pour que l'on puisse en tirer des
définitions rigoureuses.

\end{document}
